\documentclass[a4paper,12pt]{article}
\usepackage[margin=1cm]{geometry}
\usepackage[utf8]{inputenc}
\usepackage[spanish]{babel}
\usepackage{amsmath}
\usepackage{enumitem}

\pagestyle{empty}
\setlist[enumerate]{nosep, leftmargin=1.5em}

% Se cambia \hfill por un espacio simple para evitar saltos de línea
\newcommand{\prob}[2]{\item #1 \textbf{(#2 pts)}}

\begin{document}
\noindent
\begin{minipage}[t][0.47\textheight][t]{0.48\textwidth}
  \begin{center}
    \textbf{\large Mecánica de los Fluídos}\\
    \textbf{\Large Evaluación - Tema 1}
  \end{center}
  \begin{enumerate}
    \prob{Un neumático indica una presión de $2,5$ bar. Determine el valor en Pascales (Pa).}{1}
    \prob{En un sistema hidráulico se registra una presión de $30$ psi. Calcule la presión equivalente en Pa.}{1}
    \prob{La densidad de un fluido es de $1,0$ g/cm$^3$. Realice la conversión a kg/m$^3$.}{1}
    \prob{¿Cuál es la presión hidrostática en una piscina a la profundidad de $h=7,6$ m?}{1,5}
    \prob{En una prensa hidráulica, si $F_1=100$ N, $A_1=0,001$ m$^2$ y $A_2=0,1$ m$^2$, calcule la fuerza de salida $F_2$.}{1,5}
    \prob{Conociendo $F_2=100$ N, $A_1=10$ cm$^2$ y $A_2=100$ cm$^2$ de una prensa, calcule la fuerza de entrada $F_1$.}{2}
    \prob{Sobre el plato menor de una prensa se coloca una masa de $m_1=6$ kg. Si el radio del plato mayor es triple ($r_2=3r_1$), ¿qué masa se podrá levantar?}{2}
  \end{enumerate}
\end{minipage}
\hfill
\begin{minipage}[t][0.47\textheight][t]{0.48\textwidth}
  \begin{center}
    \textbf{\large Mecánica de los Fluídos}\\
    \textbf{\Large Evaluación - Tema 2}
  \end{center}
  \begin{enumerate}
    \prob{Sobre el plato menor de una prensa se coloca una masa de $m_1=8$ kg. Si el radio del plato mayor es el doble ($r_2=2r_1$), ¿qué masa se podrá levantar?}{2}
    \prob{Conociendo $F_2=200$ N, $A_1=5$ cm$^2$ y $A_2=50$ cm$^2$ de una prensa, calcule la fuerza de entrada $F_1$.}{2}
    \prob{En una prensa hidráulica, si $F_1=150$ N, $A_1=0,002$ m$^2$ y $A_2=0,2$ m$^2$, calcule la fuerza de salida $F_2$.}{1,5}
    \prob{¿Cuál es la presión hidrostática en una piscina a la profundidad de $h=12,4$ m?}{1,5}
    \prob{La densidad de un fluido es de $0,8$ g/cm$^3$. Realice la conversión a kg/m$^3$.}{1}
    \prob{En un sistema hidráulico se registra una presión de $45$ psi. Calcule la presión equivalente en Pa.}{1}
    \prob{Un neumático indica una presión de $3,2$ bar. Determine el valor en Pascales (Pa).}{1}
  \end{enumerate}
\end{minipage}

\vspace{0.2cm}

\noindent
\begin{minipage}[t][0.47\textheight][t]{0.48\textwidth}
  \begin{center}
    \textbf{\large Mecánica de los Fluídos}\\
    \textbf{\Large Evaluación - Tema 3}
  \end{center}
  \begin{enumerate}
    \prob{La densidad de un fluido es de $1,2$ g/cm$^3$. Realice la conversión a kg/m$^3$.}{1}
    \prob{En una prensa hidráulica, si $F_1=80$ N, $A_1=0,0005$ m$^2$ y $A_2=0,05$ m$^2$, calcule la fuerza de salida $F_2$.}{1,5}
    \prob{Un neumático indica una presión de $1,8$ bar. Determine el valor en Pascales (Pa).}{1}
    \prob{Sobre el plato menor de una prensa se coloca una masa de $m_1=5$ kg. Si el radio del plato mayor es cuatro veces el menor ($r_2=4r_1$), ¿qué masa se podrá levantar?}{2}
    \prob{En un sistema hidráulico se registra una presión de $22$ psi. Calcule la presión equivalente en Pa.}{1}
    \prob{¿Cuál es la presión hidrostática en una piscina a la profundidad de $h=5,2$ m?}{1,5}
    \prob{Conociendo $F_2=150$ N, $A_1=12$ cm$^2$ y $A_2=120$ cm$^2$ de una prensa, calcule la fuerza de entrada $F_1$.}{2}
  \end{enumerate}
\end{minipage}
\hfill
\begin{minipage}[t][0.47\textheight][t]{0.48\textwidth}
  \begin{center}
    \textbf{\large Mecánica de los Fluídos}\\
    \textbf{\Large Evaluación - Tema 4}
  \end{center}
  \begin{enumerate}
    \prob{Conociendo $F_2=250$ N, $A_1=8$ cm$^2$ y $A_2=80$ cm$^2$ de una prensa, calcule la fuerza de entrada $F_1$.}{2}
    \prob{¿Cuál es la presión hidrostática en una piscina a la profundidad de $h=18,1$ m?}{1,5}
    \prob{En un sistema hidráulico se registra una presión de $38$ psi. Calcule la presión equivalente en Pa.}{1}
    \prob{Un neumático indica una presión de $4,1$ bar. Determine el valor en Pascales (Pa).}{1}
    \prob{Sobre el plato menor de una prensa se coloca una masa de $m_1=10$ kg. Si el radio del plato mayor es $2,5$ veces el menor ($r_2=2,5r_1$), ¿qué masa se podrá levantar?}{2}
    \prob{La densidad de un fluido es de $0,9$ g/cm$^3$. Realice la conversión a kg/m$^3$.}{1}
    \prob{En una prensa hidráulica, si $F_1=120$ N, $A_1=0,0015$ m$^2$ y $A_2=0,15$ m$^2$, calcule la fuerza de salida $F_2$.}{1,5}
  \end{enumerate}
\end{minipage}

\end{document}